%%%%%%%%%%%%%%%%%%%%%%%%%%%%%%%%%%%%%%%%%%%%%%%%%%%%%%%%%%%%
% 文件名：卟学墨环与中微子手性_函子推导_v11.tex
% 编译：xelatex
% 编码：UTF-8
% 定位：卟学框架内的形式化表述——函子推导，非标准物理论文
% v11：定义1.1开头补公理 I 声明（点胚内部 0=1），
%      结论第2条锚点自足；其余逐字保留 v10
%%%%%%%%%%%%%%%%%%%%%%%%%%%%%%%%%%%%%%%%%%%%%%%%%%%%%%%%%%%%
\documentclass[11pt]{ctexart}

\usepackage{amsmath,amssymb,amsthm,mathrsfs}
\usepackage{geometry}
\geometry{a4paper,margin=2.5cm}
\usepackage{hyperref}

\theoremstyle{plain}
\newtheorem{thm}{定理}[section]
\newtheorem{prop}[thm]{命题}
\newtheorem{cor}[thm]{推论}

\theoremstyle{definition}
\newtheorem{defn}[thm]{定义}
\newtheorem{rem}[thm]{注记}

\newcommand{\Cat}{\mathbf}

\begin{document}

\title{墨环单向性与中微子左旋的函子对应\\
{\large ——卟学形式化接口 v11：复指数延拓}}
\author{作者}
\date{2026年10月14日}
\maketitle

\begin{abstract}
本文在 v8 手性函子推导的基础上，将源范畴的点胚从实数指数 $0^{p/q}$
升级为复数指数 $0^{z}$（$z\in\mathbb{C}$）。核心结论：
\textbf{实情形}（$z=p/q\in\mathbb{R}$）逐字继承 v8——中微子左旋是函子
$F^{p/q}$ 的必然像，$q=2$ 回墨环，$q=3$ 给三代味结构；
\textbf{复情形}（$z=a+ib$，$b\neq 0$）在同一函子框架下给出振荡相位与
CP 破坏的结构接口。本文明确区分三个层次：已证（实情形）、结构已定
（复情形）、接口开放（定量细节），为卟学框架的后续延拓保留入口。
\end{abstract}

\section{源范畴：复指数点胚}

\begin{defn}[点胚 $0^z$]\label{def:pointseed}
卟学两条原始公理提供本定义的基础：
\textbf{公理 I}：点胚内部 $0=1$（内部同一性，非对称性来源）；
\textbf{公理 II}：$1/0=\infty$（形式延拓基础）。

以下 $0^z := e^{z\ln 0}$ 是\textbf{卟学框架内的形式延拓}，
不要求经典复分析中的收敛性；$\ln 0 = -\infty$ 由公理 II
提供，作为形式符号使用，不视为经典复分析定理。
对 $z = a+ib$（$a,b\in\mathbb{R}$）按经典延拓分解：
\[
0^{a+ib} = 0^a\cdot 0^{ib},\qquad 0^{ib} = e^{-ib\infty}.
\]
$0^{ib}$ 在 $b\neq 0$ 时是一个形式相位符号，不作为经典复分析对象；
其物理尺度由延拓接口 \S\ref{sec:extension} 给出。
\end{defn}

\begin{rem}[实部与虚部的地位]\label{rem:ab}
$a$ 与 $b$ 是同一个复数 $z$ 的实部与虚部，二者地位平等，无主次之分。
在函子框架中：
\begin{itemize}
\item 实部 $a$：给出 $q$ 分支结构（$a = p/q$ 有理时），经函子映为味结构；
\item 虚部 $b$：给出相位因子 $e^{-ib\infty}$，经函子映为振荡相位。
\end{itemize}
两者的分工不是外加约定，而是复结构在函子下的两种读出。
\end{rem}

\begin{defn}[源范畴 $\Cat{Pu}_z$]\label{def:puz}
对象：由复指数 $0^z$ 生成的几何结构。
\begin{itemize}
\item $D^z$：未扭复带（$z$ 决定分支数与相位结构，局部两面/全局两面）
\item $M^z$：已扭复带（同上，局部两面/全局一面）
\end{itemize}
生成态射：
\[
t^z: D^z \to M^z
\]
（在每个分支上执行一次扭结，保持 $\mathbb{Z}_q$ 循环对称与相位结构）。
恒等态射与复合规则与 v8 的 $\Cat{Pu}$ 一致。
$\Cat{Pu}_z$ 是由对象 $\{D^z,M^z\}$ 与生成态射 $\{t^z\}$ 生成的自由范畴。
\end{defn}

\begin{rem}[$z\in\mathbb{R}$ 时回到 v8]\label{rem:backtov8}
$z = p/q\in\mathbb{R}$（$b=0$）时：
\begin{itemize}
\item $q=2$：$0^{1/2}$ 的单值化是墨环，$\Cat{Pu}_{1/2} = \Cat{Pu}$（v8 源范畴）；
\item $q=3$：$0^{1/3}$ 给出三分支结构，对应三代味；
\item 一般 $q$：$0^{1/q}$ 给 $q$ 分支。
\end{itemize}
\end{rem}

\begin{rem}[$z\in\mathbb{C}$ 时新增相位结构]\label{rem:complexcase}
$z = a+ib$（$b\neq 0$）时：
\begin{itemize}
\item $a = p/q$ 给 $q$ 分支结构（与实情形相同）；
\item $b$ 给相位 $e^{-ib\infty}$，在 $\Cat{Pu}_z$ 中表现为复带上的 $\mathrm{U}(1)$ 联络。
\end{itemize}
\end{rem}

\section{靶范畴：费米子与弱作用}

\begin{defn}[范畴 $\Cat{Ferm}$]\label{def:ferm}
对象：费米子旋量结构。
\begin{itemize}
\item $\Psi$：弱作用顶点作用下的初始态，具有左右手性分量的费米子场
（双态，局部双态/全局双态）
\item $\Psi_L$：参与弱作用后的费米子场（仅左旋分量存活，单态，
局部双态/全局单态）
\end{itemize}
态射：手性操作。
\begin{itemize}
\item $w: \Psi \to \Psi_L$：弱作用顶点（将双态费米子投影为单态手性费米子）
\end{itemize}
恒等态射与复合规则与 v8 一致。投影要求：
\begin{enumerate}
\item （幂等性）存在包含态射 $\iota: \Psi_L \hookrightarrow \Psi$ 使得
$w\circ\iota = \mathrm{id}_{\Psi_L}$。
\item （完全性）不动点集 $\mathrm{Fix}(w) = \Psi_L$，
即 $\Psi\setminus\Psi_L$ 中没有不动点。
\end{enumerate}
\end{defn}

\begin{rem}[物理对应]\label{rem:physical}
$\Psi$ 对应标准模型中参与 $SU(2)_L$ 弱作用的费米子场在弱作用顶点
作用下的初始态（如轻子二重态 $(\nu_L,e_L)$ 与右旋伴场 $e_R$ 共存的描述）。
$\Psi_L$ 对应弱作用只耦合的左手征分量。
$w$ 对应弱作用顶点处的手性投影。
\end{rem}

\begin{rem}[靶唯一性是定义而非定理]\label{rem:targetbydef}
在 $\Cat{Ferm}$ 中，从 $\Psi$ 出发的非恒等态射按定义只有 $w$ 一个，
故 $w$ 的靶 $\Psi_L$ 由定义直接给定。弱作用顶点只有一种
（$SU(2)_L$ 手性投影）是最小标准模型的物理输入，不是 $\Cat{Ferm}$
的数学推论。
\end{rem}

\begin{rem}[未走过扭结的费米子]\label{rem:untwisted}
右旋电子 $e_R$ 不参与 $SU(2)_L$ 弱作用：它只携带 $U(1)_Y$ 超荷
（$Y=-2$，本文采用 $Q = T_3 + Y/2$ 约定），通过电弱混合后对 $Z$
玻色子有部分耦合，但不进入 $SU(2)_L$ 二重态。这对应 $\Cat{Pu}_z$
中未经历扭结的复带 $D^z$。若引入 sterile $\nu_R$，同理——不参与
弱作用，不经历手性投影。
\end{rem}

\section{函子}

\begin{defn}[复函子 $F^z: \Cat{Pu}_z \to \Cat{Ferm}$]\label{def:functor}
对象映射：
\begin{align*}
F^z(D^z) &= \Psi \\
F^z(M^z) &= \Psi_L
\end{align*}
态射映射：
\begin{align*}
F^z(t^z) &= w \\
F^z(\mathrm{id}_{D^z}) &= \mathrm{id}_\Psi \\
F^z(\mathrm{id}_{M^z}) &= \mathrm{id}_{\Psi_L}
\end{align*}
复合保持与 v8 一致：$F^z(t^z\circ\mathrm{id}_{D^z}) =
F^z(t^z)\circ F^z(\mathrm{id}_{D^z}) = w$。
\end{defn}

\begin{rem}[函子保持结构]\label{rem:preserve}
$F^z$ 将 $t^z$ 的源对象 $D^z$（两面结构）映为 $\Psi$（双态结构），
将靶对象 $M^z$（一面结构）映为 $\Psi_L$（单态结构）。
$z\in\mathbb{R}$ 时与 v8 的 $F$ 完全一致。
$z\in\mathbb{C}$ 时 $F^z$ 额外将复带的 $\mathrm{U}(1)$ 结构映为
$\Psi_L$ 上的相位结构（见 \S\ref{sec:osc}）。
\end{rem}

\begin{rem}[$z=p/q$ 时逐字继承 v8]\label{rem:inherit}
$z = p/q\in\mathbb{R}$ 时 $F^{p/q} = F$（v8 的函子），
$q=2$ 回 v8 的墨环，$q=3$ 给三代味结构。
\end{rem}

\section{左旋方向的锁定}

\begin{defn}[点胚分量与 $SU(2)_L$ 弱同位旋的对应]\label{def:su2}
点胚的两分量（$0^{p/q}$ 在 $q=2$ 时的实情形）与 $SU(2)_L$ 二重态的
弱同位旋第三分量 $T_3$ 对应如下：

第一分量对应 $T_3 = +1/2$ 的态（$\nu_L$，由电荷公式 $Q = T_3 + Y/2$
且 $Y=-1$ 得 $Q=0$）；
第二分量对应 $T_3 = -1/2$ 的态（$e_L$，得 $Q=-1$）。

轻子二重态基为 $\begin{pmatrix}\nu_L\\ e_L\end{pmatrix}$，
上分量 $\nu_L$ 对应 $T_3=+1/2$，下分量 $e_L$ 对应 $T_3=-1/2$。

（注：$\nu_L$ 位于二重态上分量并带 $T_3=+1/2$，是标准模型的约定；
若换用相反排序，本定义的对应需同步翻转。）
\end{defn}

\begin{rem}[手性与同位旋的独立性]\label{rem:independence}
定义~\ref{def:su2} 建立了点胚分量与 $T_3$ 的对应，但这是卟学点胚
与 $SU(2)_L$ 二重态之间的\emph{对应约定}，而非从卟学内部推导出来的
结论。卟学只给出点胚有两个分量；至于为什么第一分量对应 $T_3=+1/2$
而非 $-1/2$，这是将一个二元对应选定为 $0\leftrightarrow +1/2$、
$1\leftrightarrow -1/2$——换一种选定在卟学内部同样自洽。

此外，$T_3$ 与手性（$P_L/P_R$）是两个独立的量子数。
$SU(2)_L$ 的表示结构本身不区分左右手性——它只规定二重态的 $T_3$ 值。
“轻子二重态的分量是左手征的”这一事实来自最小标准模型的实验输入。

因此函子 $F^z$ 将扭结方向性映为 $P_L$ 的锁定机制是\emph{两层结构
共同作用的}：
\begin{enumerate}
\item 对应约定层（卟学侧）：点胚分量与 $SU(2)_L$ 二重态 $T_3$ 的对应；
\item 经验层（物理侧）：最小标准模型中弱作用只耦合左手征。
\end{enumerate}
两层结合，函子方向被确定。
\end{rem}

\section{手性定理（实情形）}

\begin{thm}[给定最小标准模型结构下的函子必然性]\label{thm:main}
设 $z = p/q\in\mathbb{R}$。在 $\Cat{Pu}_{p/q}$ 中，点胚 $0^{p/q}$ 的
非对称性生成扭结态射 $t^{p/q}: D^{p/q}\to M^{p/q}$，
这是自由生成范畴中从“两边”到“一边”的唯一生成态射。

由函子 $F^{p/q}$（定义~\ref{def:functor}），$F^{p/q}(t^{p/q}) = w$
是 $\Cat{Ferm}$ 中对应的“双态$\to$单态”态射，
其靶按定义给定为 $\Psi_L$。

由定义~\ref{def:su2} 和注记~\ref{rem:independence}，
卟学---$SU(2)_L$ 对应约定与最小标准模型经验事实共同锁定方向为
左旋投影 $P_L$。

因此在给定函子 $F^{p/q}$ 与最小标准模型结构下，中微子左旋是
墨环扭结结构通过函子强制传输到费米子范畴的必然结果。
\end{thm}

\begin{proof}
与 v8 定理 6.1 的证明逐字一致：
点胚 $0^{p/q}$ 的非对称性在自由范畴 $\Cat{Pu}_{p/q}$ 中生成唯一态射
$t^{p/q}: D^{p/q}\to M^{p/q}$；
函子 $F^{p/q}$ 保持此结构，$F^{p/q}(t^{p/q}) = w$ 将双态对象 $\Psi$
映为单态对象 $\Psi_L$，靶由注记~\ref{rem:targetbydef} 按定义给定；
方向性由两层锁定（定义~\ref{def:su2} 与注记~\ref{rem:independence}），
确定 $w = P_L$；
在最小标准模型中，中微子只以左手征分量进入轻子二重态
$(\nu_L,e_L)$，故 $\Psi_L$ 在轻子扇区的实例正是 $\nu_L$。
\end{proof}

\begin{rem}[$q=2$ 回 v8，$q=3$ 给三代味]\label{rem:qcases}
$p/q = 1/2$ 时定理~\ref{thm:main} 即 v8 的中微子左旋定理。
$p/q = 1/3$ 时 $\Psi_L$ 自动获得 $\mathbb{Z}_3$ 味结构，
三代味来自 $0^{1/3}$ 的三分支，见 \S\ref{sec:flavor}。
\end{rem}

\section{味结构：实部分支}\label{sec:flavor}

\begin{rem}[味从 $a$ 的分支结构来]\label{rem:flavor}
$z$ 的实部 $a = p/q$ 给出 $q$ 分支结构，经 $F^{p/q}$ 映为 $\Psi_L$
的 $\mathbb{Z}_q$ 味结构：
\[
\Psi_L = \bigoplus_{k=0}^{q-1} \Psi_L^{(k)}.
\]
味本征态不是分支的不动点，而是 $\mathbb{Z}_q$ 分支置换的本征态。
$q=3$ 给出三代味。
\end{rem}

\begin{rem}[味结构是函子像的内蕴属性]\label{rem:flavorfrom}
$\Psi_L$ 的 $\mathbb{Z}_q$ 味结构不是外加输入，
它是 $F^{p/q}$ 从 $M^{p/q}$ 的 $q$ 分支结构直接传输的像。
$\Cat{Pu}_{p/q}$ 的每个分支 $\leftrightarrow$ $\Psi_L$ 的每个味分量。
\end{rem}

\section{振荡与 CP：虚部相位}\label{sec:osc}

\begin{rem}[振荡：$b$ 的相位结构]\label{rem:osc}
$z = a+ib$ 的虚部 $b$ 经 $F^z$ 映为相位因子
\[
e^{-ib\infty},
\]
在 $\Cat{Ferm}$ 中表现为 $\Psi_L$ 的复相位。
在味基下，振荡态射具有形式
\[
\mathrm{osc}: \Psi_L^{(k)} \longmapsto \sum_{l=0}^{q-1}
U_{kl}\, e^{i\theta_l}\,\Psi_L^{(l)},
\]
其中 $U_{kl}$ 的\textbf{结构来源}是 $\mathbb{Z}_q$ 分支置换，
$e^{i\theta_l}$ 是分支 $l$ 的相位累积（$b$ 的像）。

\textbf{嵌入说明。} $\mathbb{Z}_q$ 的正则表示是 $q\times q$
循环置换矩阵，它是 $\mathrm{U}(q)$ 的离散子群，不包含连续的混合角。
从 $\mathbb{Z}_q$ 分支置换到物理上 $SU(q)$ 混合矩阵
（$q=3$ 时为 PMNS 的 $3\times 3$ 幺正矩阵）的精确嵌入
\textbf{留待延拓}（见 \S\ref{sec:extension}）；当前版本只给出
$U_{kl}$ 的结构来源是分支置换，不声称已给出 $SU(q)$ 的连续参数。

PMNS 混合角与振荡频率的定量对应留待延拓。
\end{rem}

\begin{rem}[CP 破坏：复指数的非实性]\label{rem:cp}
CP 变换在源侧对应 $z\mapsto\bar z$（复共轭）。
$b\neq 0$ 时 $0^z\neq 0^{\bar z}$，即复指数的非实性。

\textbf{函子桥接。} 在靶侧，CP 变换对应 PMNS 矩阵的复共轭
$U_{PMNS}\mapsto U_{PMNS}^*$。$b\neq 0$ 时
$F^z\neq F^{\bar z}$，故经函子传输后
$U_{PMNS}$ 不能被实正交矩阵吸收——
其复相位不可消除，即 CP 破坏。

$q\geq 3$ 时 $\mathbb{Z}_q$ 的复特征 $\omega^k$ 与 $\omega^{-k}$ 不等，
与 CP 破坏相位不可消除一致。
$q = 2$ 时 $\mathbb{Z}_2$ 自共轭，$U_{PMNS}$ 可被实正交矩阵吸收，
无 CP 破坏。定量形式（PMNS 相位角的精确值）留待延拓。
\end{rem}

\begin{rem}[$q\geq 3$ 与 CP 破坏的物理一致性]\label{rem:km}
物理上 CP 破坏在标准模型中要求至少三代（Kobayashi-Maskawa）。
本框架中 $q\geq 3$ 时 $\mathbb{Z}_q$ 分支结构非自共轭，
与物理条件结构性一致。
\end{rem}

\section{延拓接口}\label{sec:extension}

\begin{rem}[当前版本的三层边界]\label{rem:boundary}
本版本（v11）在以下意义上完成：
\begin{enumerate}
\item \textbf{已证}（$z\in\mathbb{R}$）：实情形全部定理逐字成立
（\S\ref{sec:flavor} 之前的全部内容，以及定理~\ref{thm:main}）。
\item \textbf{结构已定}（$z\in\mathbb{C}$，$b\neq 0$）：
$0^z = 0^a\cdot 0^{ib}$ 给出实部/虚部两层结构；
实部给分支（味），虚部给相位（振荡），
经函子 $F^z$ 分别映为 $\Psi_L$ 的味结构与相位结构。
\item \textbf{接口开放}（定量细节）：
  \begin{itemize}
  \item \textbf{$\mathbb{Z}_q$ 到 $SU(q)$ 的嵌入}：
  从分支置换到 PMNS 连续混合矩阵的精确嵌入方式
  （$q=3$ 时给出三个混合角与一个 CP 相位）；
  \item $b$ 的物理尺度：振荡频率 $\theta_l$ 与质量平方差
  $\Delta m^2$、基线 $L$、能量 $E$ 的定量对应；
  \item $a$ 的有理逼近规则：分支数 $q$ 如何在 $a$ 连续时涌现；
  \item $\infty$ 在指数映射下的周期结构：$e^{-ib\infty}$ 的
  $2\pi$ 周期性从公理 III 如何精确导出。
  \end{itemize}
\end{enumerate}
这些不是 v11 的缺失，而是卟学框架作为可延拓体系的设计：
当前版本给出结构和已证部分，细节由后续延拓按物理对接补全。
\end{rem}

\begin{rem}[方法论说明]\label{rem:method}
本版本延续 v8 的处理方式：明确标注哪些是已证的、哪些是约定的、
哪些是留给延拓的。v8 标注了“$T_3$ 对应是约定”；
v11 在此基础上标注“$b$ 的尺度是接口”与“$\mathbb{Z}_q\to SU(q)$
嵌入是接口”。
这一姿态与卟学“自然本无禁区”的框架一致：
不因为存在开放接口而拒绝框架，也不把未证明的东西当作已证明。
\end{rem}

\section{结论}

\begin{enumerate}
\item 三代味从 $0^{1/3}$ 的实分支来（$a = 1/3$）。
\item 中微子左旋从 $0^{p/q}$ 每分支的内部同一性（公理 I：点胚内
$0=1$）来（实情形，定理~\ref{thm:main}）。
\item 振荡与 CP 从 $0^z$ 的复结构来（$b\neq 0$），
经函子 $F^z$ 映为 $\Psi_L$ 的相位结构；定量细节留待延拓。
\item 全部结构来自 $0^z = e^{z\ln 0}$ 的经典延拓。
\item 实情形封闭（v8 全部定理逐字继承），
复情形接口开放（\S\ref{sec:extension}），
与卟学“自然本无禁区”的框架一致。
\end{enumerate}

\appendix
\section{附录：态射与对象的转化}

在范畴论中，态射 $f: A\to B$ 可以通过以下方式转化为对象：
\begin{enumerate}
\item \textbf{箭头范畴} $\mathrm{Arr}(\mathcal{C})$：其对象为
$\mathcal{C}$ 中的态射 $f: A\to B$，态射为交换正方形。
在 $\mathrm{Arr}(\Cat{Pu}_z)$ 中，$t^z$ 本身是一个对象。
\item \textbf{像对象}：$\mathrm{im}(t^z) = M^z$ 作为 $\Cat{Pu}_z$ 中的对象。
\item \textbf{切片范畴} $\mathcal{C}/M^z$：对象为所有以 $M^z$ 为靶的态射。
\end{enumerate}

函子 $F^z$ 在 $\Cat{Pu}_z$ 中作为普通函子定义；
$F^z$ 诱导函子 $\mathrm{Arr}(F^z): \mathrm{Arr}(\Cat{Pu}_z)
\to \mathrm{Arr}(\Cat{Ferm})$，将对象 $t^z$ 映为对象 $w$。

\section{附录：v8 到 v11 的对应表}

\begin{center}
\begin{tabular}{lll}
\hline
\textbf{结构} & \textbf{v8} & \textbf{v11} \\
\hline
点胚 & $\vartheta = (0,1)$ & $0^z$，$z\in\mathbb{C}$ \\
源范畴 & $\Cat{Pu}$（$q=2$） & $\Cat{Pu}_z$（$q$ 由 $a$ 决定） \\
生成态射 & $t: D\to M$ & $t^z: D^z\to M^z$ \\
靶范畴 & $\Cat{Ferm}$ & $\Cat{Ferm}$（不变） \\
函子 & $F$ & $F^z$ \\
味结构 & （无） & $\Psi_L = \bigoplus_{k=0}^{q-1}\Psi_L^{(k)}$ \\
相位结构 & （无） & $e^{-ib\infty}$ \\
手性定理 & 定理 6.1 & 定理~\ref{thm:main}（实情形） \\
\hline
\end{tabular}
\end{center}

$z = 1/2$（$b=0$）时 v11 与 v8 完全一致。

\end{document}